\section{Joint sampling of leaders and removals}
\label{sam:section}

This section constructs the prefixes used by the recovery algorithm.  Its
input is a real solution $S$, an upper bound $M_0$ on the number of bad
comparison centers, and guesses for two integers $P$ and $m$.  The comparison
solution, its partition, and its proxies are used only in the analysis.
For each guess, the algorithm samples a client proportionally to its current
reference distance. A fair coin selects an attempt to add a guessed ball
centered at that client or to delete its assigned original center.
Every valid prefix is saved for completion; the comparison-dependent
stopping rules below certify that one such prefix is useful.

Distance-proportional sampling has a general clustering antecedent in
\cite[Section~5]{ArthurVassilvitskii2007}.  The closer predecessor here
is the recovery construction of \cite[Sections~4.2--4.4]{CGLSS2026},
which samples leaders, guesses balls and deletes original centers.
Our experiment interleaves these operations.  The population radii and
their amortized menu cost are what make the dependence on the number
of bad clusters single-exponential.

All distances in the instance considered here are positive integers between
distinct indexed points, with a polynomial upper bound.  The reduction to
this case is proved in the preliminaries.

Fix, for the analysis, a comparison solution $O$ with
$|O|=|S|=h$ and $\operatorname{cost}(O)=P>0$.  Fix nearest assignments to $O$,
write $C_i$ for the cluster of $i\in O$, and put
\[
 n_i=|C_i|,\qquad P_i=\sum_{p\in C_i}d_p^*,\qquad
 d_p^*=d(p,O),\qquad
 \bar d_i=P_i/n_i\quad(n_i>0).
\]
Give an empty cluster mean zero.  Let $\mathcal G$ and $\mathcal B$ be the
good and bad centers, respectively, and suppose that $m=|\mathcal B|$ and
there are distinct proxies $\phi_i\in S$, $i\in\mathcal G$, such that
\begin{equation}
 \sum_{i\in\mathcal G}\sum_{p\in C_i}d(p,\phi_i)
 \le P_g+0.01P,
 \qquad P_g=\sum_{i\in\mathcal G}P_i.
 \label{sam:proxy-hypothesis}
\end{equation}
Thus
\begin{equation}
 S_0=S\setminus\{\phi_i:i\in\mathcal G\},\qquad |S_0|=m.
 \label{sam:nonproxy-set}
\end{equation}
The set $S_0$ need not be found.  If $m=0$, the proxies exhaust $S$ and
Equation~\eqref{sam:proxy-hypothesis} already gives
$\operatorname{cost}(S)\le1.01P$.  In what follows $m\ge1$.
The main recovery algorithm uses $\kappa=10^{-6}$.  Some geometric and
probabilistic lemmas below allow any fixed rational $0<\kappa\le1$ so that
they can also be used for the accuracy refinement.

\subsection{Computable radii and reference metrics}

For a client $w$ and $q\ge0$, let
$\mathcal D(w,q)=\{p\in D:d(w,p)\le q\}$ be the closed client ball.
Define
\begin{equation}
 r^0(w)=\min\{q\ge0:q|\mathcal D(w,q)|\ge P/m\}.
 \label{sam:r0-definition}
\end{equation}
This radius depends only on the guessed $P,m$ and the input metric.

\begin{lemma}[The computable radius]\label{sam:computable-radius}
If $d_1\le\cdots\le d_n$ are the client distances from $w$, including the
zero distance from $w$ to itself, then
\begin{equation}
 r^0(w)=\min_{1\le j\le n}\max\{d_j,P/(mj)\}>0.
 \label{sam:r0-formula}
\end{equation}
In particular the minimum in Equation~\eqref{sam:r0-definition} is
attained and its closed ball satisfies the defining population inequality.
Suppose that $A\subseteq\mathcal B$ and each $i\in A$ has a client
$w_i\in C_i$ with $r_i=d(w_i,i)\le a\bar d_i$, where $a>0$.  Then
\begin{equation}
 \begin{aligned}
 r^0(w_i)&\le(a+2)\bar d_i+\frac{2P}{mn_i},\\
 \sum_{i\in A}n_i r^0(w_i)
 &\le(a+2)\sum_{i\in A}P_i+\frac{2|A|}{m}P
 \le(a+4)P.
 \end{aligned}
 \label{sam:leader-radius-sum}
\end{equation}
The last bound is $5.2P$ when $a=1.2$, and is at most $5.5P$ when
$a\le3/2$.
\end{lemma}

\begin{proof}
For every $j$, the number $q_j=\max\{d_j,P/(mj)\}$ has at least $j$
clients in its closed ball and $q_jj\ge P/m$, so it is feasible.  Conversely,
if $q$ is feasible and $j=|\mathcal D(w,q)|$, then $j\ge1$,
$d_j\le q$, and $P/(mj)\le q$.  Hence $\min_jq_j\le q$.  These two
observations prove Equation~\eqref{sam:r0-formula}, attainment, and
positivity.

For a selected leader the cluster is nonempty.  At least $n_i/2$ of its
clients have $d_p^*\le2\bar d_i$: otherwise their distances alone would
sum to more than $P_i$.  When $\bar d_i=0$ all clients have zero distance,
so the same assertion holds.  These clients are within
$(a+2)\bar d_i$ of $w_i$.  At radius
$(a+2)\bar d_i+2P/(mn_i)$ their population certifies feasibility in
Equation~\eqref{sam:r0-definition}.  Multiply the resulting bound by $n_i$
and sum.  Distinct leaders here belong to distinct bad clusters, so
$|A|\le m$ and $\sum_{i\in A}P_i\le P$.
\end{proof}

Let $\Delta>0$ be the diameter of the original metric on $D\cup F$.
Adjoin a permanent point $\delta_\infty$ at distance $\Delta$ from every
original point.  For each selected pair $(w_i,R_i)$ with $R_i>0$, adjoin
a distinct leaf $\delta_i$ at $w_i$ with edge length $R_i$.
This leaf construction is used in \cite[Section~4.3.1]{CGLSS2026};
we record its metric and tie-breaking properties for the present process.
Write
$\Lambda=\{\delta_i:i\in A\}$ and
\begin{equation}
 H=(S\setminus Q)\cup\Lambda\cup\{\delta_\infty\},\qquad
 \mathcal F_i=\{x\in F:d(w_i,x)\le R_i\}.
 \label{sam:reference}
\end{equation}
Here $Q$ is the set of original centers removed so far.  The labels $i\in A$
describe a correct prefix; the algorithm stores only the list of leaders,
radii, and leaves.

\begin{lemma}[Dummy distances and ties]\label{sam:dummy-metric}
The construction above extends the original metric without changing its
distances.  For every original point $x$,
\[
 d(x,\delta_i)=d(x,w_i)+R_i,\qquad
 d(\delta_\infty,\delta_i)=\Delta+R_i.
\]
If $R_i\ge d(w_i,O)$, then $d(x,\delta_i)\ge d(x,O)$ for every original
point $x$.  If $x_i\in\mathcal F_i$, then
$d(x,x_i)\le d(x,\delta_i)$ for every original point $x$.
Every nonempty real solution similarly dominates the permanent dummy.

There is a fixed assignment rule, independent of $O,S_0$, and $Q$, such
that after any state in which further operations only add leaves, the
client set assigned to any fixed subset of surviving original centers can
only shrink.
\end{lemma}

\begin{proof}
One may first form the complete graph on the original metric points with
their metric edge lengths and join $\delta_\infty$ to all of them by edges
of length $\Delta$.  A path through $\delta_\infty$ cannot shorten an
original distance, since its two adjacent edges have total length
$2\Delta$.  Taking the resulting shortest-path metric and then attaching
leaves proves the distance formulas and all triangle inequalities.  In
particular, for different leaves the distance is
$R_i+d(w_i,w_j)+R_j$.

The distance-to-set function is $1$-Lipschitz.  More directly,
$d(x,O)\le d(x,w_i)+d(w_i,O)\le d(x,\delta_i)$, whereas
$d(x,x_i)\le d(x,w_i)+d(w_i,x_i)\le d(x,\delta_i)$.
Every original facility is within $\Delta$ of each original client, which
proves the permanent-dummy assertion for a nonempty solution.

Fix an order on the original facilities, put the permanent dummy after
them, and order leaves by creation time.  Assign each client to the first
nearest available reference point in this order.  Leaf addition preserves
the distances and relative order of all existing choices.  A client whose
chosen reference point was outside a given fixed subset of original
centers cannot begin to choose a point in that subset after more choices
are added.  This proves the assertion simultaneously for every such
subset, including a subset unknown to the algorithm.
\end{proof}

All actual assignments below use this rule.  Write $J_s$ for the clients
assigned to an original center $s\in S\setminus Q$, and put
\[
 W(H)=\sum_{p\in D}d(p,H).
\]
The reference is always nonempty because of $\delta_\infty$.

\subsection{A finite radius experiment}

We first specify a radius experiment that does not test any condition
involving the unknown comparison solution.  Given $w$, put $q=r^0(w)$ and
define the integers
\[
 \begin{aligned}
 M_\kappa&=\lceil10/\kappa\rceil+1,&
 J_{\kappa,m}&=1+\lceil2m/\kappa\rceil,\\
 E_\kappa&=\min\{e\in\mathbb Z_{\ge0}:(1+\kappa)^e\ge10\}.
 \end{aligned}
\]
Choose one of the following three menus with probability $1/3$ each.
\begin{enumerate}[label=\textup{(\arabic*)}]
\item Choose $j$ uniformly from $\{1,\ldots,M_\kappa\}$ and set
      $R=j\kappa q$.
\item If $\Lambda\ne\varnothing$, let $D_\Lambda=d(w,\Lambda)$,
      choose $e$ uniformly from $\{0,\ldots,E_\kappa\}$, and set
      $R=D_\Lambda/(1+\kappa)^e$.  If $\Lambda=\varnothing$, this menu
      is undefined and ends the run.
\item Choose $j\in\{1,\ldots,J_{\kappa,m}\}$ with probability
      \[
       \frac{2^{-j}}{\sum_{u=1}^{J_{\kappa,m}}2^{-u}}
       =\frac{2^{J_{\kappa,m}-j}}{2^{J_{\kappa,m}}-1},
      \]
      and set $R=j\kappa q$.
\end{enumerate}
All radii in these menus are positive.  The algorithm may discard a branch
if its facility ball is empty.  The first two menus have a number of entries
depending only on the fixed $\kappa$; the third has $O_\kappa(m)$ entries.

For the analysis, call a radius \emph{correct} at $w$ if, for
$r=d(w,O)$,
\begin{equation}
 r\le R\le(1+\kappa)r+\kappa r^0(w).
 \label{sam:upward-radius}
\end{equation}
On any history with correct previous radii, Lemma~\ref{sam:dummy-metric}
gives $D_\Lambda\ge r$ whenever $\Lambda$ is nonempty.  Set
$D_\Lambda=+\infty$ when it is empty, solely for the following analysis.
We prescribe a correct menu and entry as follows:
\begin{itemize}
\item If $r\le10q$, use menu \textup{(1)} with
      $j=\max\{1,\lceil r/(\kappa q)\rceil\}$.
\item If $r>10q$ and $D_\Lambda\le10r$, use menu \textup{(2)} with the
      largest $e\le E_\kappa$ for which
      $D_\Lambda/(1+\kappa)^e\ge r$.
\item If $r>10q$ and $D_\Lambda>10r$, use menu \textup{(3)} with
      $j=\lceil r/(\kappa q)\rceil$.
\end{itemize}
The first entry belongs to its stated range and satisfies
$r\le R\le r+\kappa q$, including $r=0$.  In the second case
$r\le D_\Lambda\le10r$, so the prescribed exponent exists and gives
$r\le R\le(1+\kappa)r$.  The next lemma proves that all prescribed
third-menu entries also lie in their finite range.  Thus these prescriptions
are analytical choices among actual branches, not tests performed by the
algorithm.

\begin{lemma}[Amortized radius guesses]\label{sam:radius-budget}
Fix $0<\kappa\le1$.  Consider a sequence with at most $m$ ball additions,
each using the prescribed correct radius above, and suppose that every
previously added leaf remains present at every later addition.  Arbitrary
removals of original centers may occur between additions.  Among additions
using the third prescribed menu, the closed client balls
$\mathcal D(w,r^0(w))$ are pairwise disjoint.  Writing $r=d(w,O)$ and
$q=r^0(w)$ for these additions, one has
\begin{equation}
 \sum_{\text{third menu}}\frac rq\le\frac{10m}{9},
 \qquad
 \sum_{\text{third menu}}j
 \le\left(\frac{10}{9\kappa}+1\right)m.
 \label{sam:third-menu-budget}
\end{equation}
Consequently each prescribed third-menu index is at most
$J_{\kappa,m}$.  There is a constant $c_\kappa\in(0,1]$ such that the
product of the probabilities of all prescribed menu-and-entry choices is
at least
\begin{equation}
 c_\kappa^m\,
 2^{-(10/(9\kappa)+1)m}=\exp(-O_\kappa(m)).
 \label{sam:menu-probability}
\end{equation}
These statements do not require a fixed order of additions or a fixed
rule for selecting their leaders.
\end{lemma}

\begin{proof}
Consider two third-menu additions, with the first earlier than the second,
and give their quantities subscripts $a,b$.  Suppose a client belongs to
both closed client balls.  Then
\[
 d(w_a,w_b)\le q_a+q_b<\frac{r_a+r_b}{10}.
\]
Both inequalities $q_a<r_a/10$ and $q_b<r_b/10$ are strict because these
are third-menu additions.  Lipschitzness of distance to $O$ gives
$|r_a-r_b|<(r_a+r_b)/10$, and hence
$9/11<r_a/r_b<11/9$.  The earlier persistent leaf is an available
competitor in $D_\Lambda$ at the later step.  Its correct radius therefore
gives
\begin{align*}
 D_\Lambda
 &\le d(w_b,w_a)+R_a\\
 &<\frac{r_a+r_b}{10}+(1+\kappa)r_a+\kappa r_a/10\\
 &\le\left(\frac1{10}+
       \frac{11}{9}\left(\frac{11}{10}+\frac{11\kappa}{10}\right)
       \right)r_b
 <3r_b<10r_b.
\end{align*}
Here the penultimate strict inequality uses $\kappa\le1$.
This contradicts the condition for prescribing the third menu at $b$.
Thus the client sets are disjoint; no ambient space containing whole
geometric balls is needed.

For a third-menu addition every $p\in\mathcal D(w,q)$ satisfies
$d(p,O)\ge r-d(w,p)\ge r-q>0.9r$.  The attained population inequality
from Lemma~\ref{sam:computable-radius} gives
$|\mathcal D(w,q)|\ge P/(mq)$.  Disjointness now implies
\[
 P=\sum_{p\in D}d(p,O)
 \ge0.9\sum_{\text{third menu}}r\,|\mathcal D(w,q)|
 \ge\frac{0.9P}{m}\sum_{\text{third menu}}\frac rq.
\]
Cancel $P>0$ to obtain the first bound in
Equation~\eqref{sam:third-menu-budget}.  For every third-menu entry,
\[
 j=\lceil r/(\kappa q)\rceil\le r/(\kappa q)+1.
\]
There are at most $m$ additions, so the second bound follows.
In particular,
\[
 j\le\frac{10m}{9\kappa}+1\le\frac{2m}{\kappa}+1\le J_{\kappa,m}.
\]
Availability follows inductively: a proposed third-menu addition is
disjoint from previous ones using only their radii, so the same population
argument certifies its prescribed index before its radius is selected.

Set
\[
 c_\kappa=\min\left\{\frac1{3M_\kappa},
                  \frac1{3(E_\kappa+1)},\frac13\right\}.
\]
The prescribed choice in either of the first two menus has probability
at least $c_\kappa$.  In the third it has probability at least
$c_\kappa2^{-j}$, since the normalizing sum is less than one.
Multiplication and Equation~\eqref{sam:third-menu-budget} prove
Equation~\eqref{sam:menu-probability}.  If fewer than $m$ additions occur,
replacing their number by $m$ only weakens the lower bound.
\end{proof}

\subsection{The experiment and adaptive guides}

The radius estimate has removed any aspect-ratio loss from a sequence
of correct guesses.  It remains to show that the random client draws
can supply suitable leaders and removals with constant progress
probability at each step.  We first give a likelihood argument that
allows these suitable choices to depend on the whole past history.

For guessed $P,m$ the actual experiment starts with $Q=\varnothing$ and no
leaves.  It records its initial state and every subsequent valid state as
a candidate prefix.  For at most $2m$ steps it performs the following
operation.
\begin{enumerate}[label=\textup{(\arabic*)}]
\item Compute $H$, its nearest assignments, and $W(H)$.  If $W(H)=0$,
      end this run.
\item Sample $w\in D$ with probability $d(w,H)/W(H)$ and independently
      choose ball mode or removal mode, with probability $1/2$ each.
\item In ball mode, run the finite radius experiment.  If it is defined
      and $\mathcal F(w,R)=\{x\in F:d(w,x)\le R\}$ is nonempty,
      append its ball and leaf.  In removal mode, remove the reference
      point assigned to $w$ if it belongs to $S\setminus Q$, and add
      it to $Q$.  An undefined operation ends the run.
\item End the run if it would have more than $m$ balls or more than $m$
      removals; otherwise record its new state.
\end{enumerate}
Only recorded states satisfying these observable size restrictions need be
sent to completion.  Prefixes on which no feasible completion exists can
be discarded.  The experiment never tests membership in $S_0$, the
identity of a comparison cluster, or a comparison-dependent stopping
condition.

\begin{lemma}[Adaptive guided trajectories]\label{sam:probability}
Fix $O,P,m$, $0<\kappa\le1$, and a number $\rho\in(0,1]$ independent of
the instance size.  Suppose an analytical guided process uses the same
initial state and, at each of its nonterminal histories, does the following:
\begin{enumerate}[label=\textup{(\roman*)}]
\item From that history it deterministically specifies ball mode or removal
      mode and a set of eligible clients
      such that, under the actual weighted client draw and fair mode
      choice, the joint probability of this event is at least $\rho$.
\item It draws the client from the actual law conditioned on that event.
      It performs exactly the actual update for the sampled outcome.
      Thus a removal deletes the original reference center assigned to
      that client, which must be a previously unremoved member of $S_0$.
      A ball
      addition has a leader in a previously uncovered bad cluster and uses
      its prescribed correct menu and entry.
\end{enumerate}
Suppose the guide stops after at most $2m$ steps on every supported
trajectory.  Leaves are retained throughout, and there are at most $m$
ball additions and $m$ removals.  The eligible sets, modes, and stopping
rules may depend on the entire preceding augmented history and on $O$;
they may also have several successive stages, with their stage determined
by that history.  Then the actual experiment records a
terminal state of this guide with probability at least
\begin{equation}
 \rho^{2m}c_\kappa^m
       2^{-(10/(9\kappa)+1)m}
 =\exp(-O_{\kappa,\rho}(m)).
 \label{sam:guided-likelihood}
\end{equation}
The same conclusion, with a changed constant in the exponent, holds with
a polynomial-time implementation using a bounded number of random bits
per operation.
\end{lemma}

\begin{proof}
First use the exact rational distributions.  Include the sampled client,
mode, menu, and menu entry in the history; entries irrelevant to a removal
are omitted.  Fix a guide-supported nonterminal history $\gamma$.
Let $E_\gamma$ be its prescribed mode-and-client event and let
$p_\gamma\ge\rho$ be that event's probability in the actual experiment.
For a particular eligible client, the ratio of its actual mode-and-client
probability to its guided conditional probability is exactly $p_\gamma$.
For a ball addition multiply this ratio by the probability of its
prescribed menu and entry, since the guide makes that choice deterministically.
Correct upward radii ensure that the ball contains the corresponding
comparison center.  Thus no supported operation is rejected for an empty
facility ball or for exceeding the observable limits.

Along any terminal guided trajectory $\gamma$ of length at most $2m$,
multiply these conditional ratios.  Lemma~\ref{sam:radius-budget} applies
to its full adaptive sequence and bounds the product of its menu losses.
Consequently the actual cylinder probability of $\gamma$ is at least the
right-hand side of Equation~\eqref{sam:guided-likelihood} times its guided
probability.  The terminal guided histories form a prefix-free family:
once the guide's stopping rule holds, it makes no further draws.
Their corresponding cylinders in the actual history tree are therefore
disjoint.  Summing over them, and using that the guide stops with total
probability one, proves the bound.  This summation does not select one
fixed sequence of leaders; it integrates over every client history generated
by the guide.  Nor does it require the actual algorithm to know the
stopping predicate: it records all prefixes, including the terminal one
on every one of these cylinder events.

For the finite-bit assertion, the rational-arithmetic and dyadic
implementation proved below preserves the probability of every elementary
transition from below by a factor $1/2$.  At most $2m$ such factors change
Equation~\eqref{sam:guided-likelihood} by at worst $2^{-2m}$, which has
the asserted form.
\end{proof}

\subsection{Near clients and the main stopping rule}

The following elementary bound also allows a variable cutoff.
\begin{lemma}[Near-client witnesses]\label{sam:near-clients}
Let $n_i>0$, $a>1$, and
$N_i(a)=\{p\in C_i:d_p^*\le a\bar d_i\}$.  For any nonempty reference
set in a metric extension, set $T_i(a)=\sum_{p\in N_i(a)}d(p,H)$.
Then
\begin{equation}
 |N_i(a)|\ge(1-1/a)n_i,\qquad
 d(i,H)\le a\bar d_i+
              \frac{T_i(a)}{(1-1/a)n_i}.
 \label{sam:general-near-bound}
\end{equation}
For the main cutoff $N_i=N_i(1.2)$ and
$T_i=\sum_{p\in N_i}d(p,H)$ one also has
\begin{equation}
 d(i,H)\le1.02\bar d_i+\frac{60T_i}{n_i},\qquad
 d(p,H)\le1.85d_p^*+\frac{60T_i}{n_i}\quad(p\in C_i\setminus N_i).
 \label{sam:far-client-bound}
\end{equation}
\end{lemma}

\begin{proof}
If $\bar d_i>0$, fewer than $n_i/a$ clients have distance strictly greater
than $a\bar d_i$, by summing their distances.  If $\bar d_i=0$, all
clients belong to $N_i(a)$.  This proves the population bound in both
cases.  Choose a client in $N_i(a)$ minimizing its reference distance;
that distance is at most $T_i(a)/|N_i(a)|$.  The triangle inequality
through this client proves Equation~\eqref{sam:general-near-bound}.

Apply the population argument with cutoff $1.02=51/50$.  The set of
clients with $d_p^*\le1.02\bar d_i$ has at least $n_i/51\ge n_i/60$
members and is contained in $N_i$.  Its member of minimum reference
distance gives the first inequality in Equation~\eqref{sam:far-client-bound}.
For $p\notin N_i$, one has $d_p^*>1.2\bar d_i$, so
\[
 d(p,H)\le d_p^*+d(i,H)
 \le\left(1+\frac{1.02}{1.2}\right)d_p^*
       +\frac{60T_i}{n_i}
 =1.85d_p^*+\frac{60T_i}{n_i}.
\]
If the mean is zero, there are no clients outside $N_i$, and this last
assertion is vacuous.
\end{proof}

A \emph{correct main state} has $Q\subseteq S_0$, a set
$A\subseteq\mathcal B$ of distinct covered centers, and exactly one
correct ball per $i\in A$, with leader $w_i\in N_i$.
Let $B=\mathcal B\setminus A$ be the uncovered bad centers and put
\begin{equation}
 \begin{split}
 T_{\rm near}(H)&=\sum_{i\in B}\sum_{p\in N_i}d(p,H),\\
 R^*&=S_0\setminus Q,\qquad
 I(H)=\sum_{s\in R^*}\sum_{p\in J_s}d(p,H).
 \end{split}
 \label{sam:stopping-quantities}
\end{equation}
For empty clusters let $N_i=\varnothing$ and $T_i=0$.
The analytical stopping rule is
\begin{equation}
 T_{\rm near}(H)\le\kappa P,
 \qquad I(H)\le\kappa P.
 \label{sam:stopping-rule}
\end{equation}

\begin{lemma}[Sampling mass]\label{sam:sampling-mass}
At every correct main state with $\kappa=10^{-6}$,
\begin{equation}
 W(H)\le3.42P+61T_{\rm near}(H).
 \label{sam:weight-bound}
\end{equation}
If the first stopping inequality fails, the joint event of sampling a
client in $\bigcup_{i\in B}N_i$ and choosing ball mode has probability
at least
\begin{equation}
 \rho_\kappa=\frac{\kappa}{2(3.42+61\kappa)}>0.
 \label{sam:progress-constant}
\end{equation}
If the first stopping inequality holds and the second fails, the same
lower bound holds for sampling a client assigned to $R^*$ and choosing
removal mode.
\end{lemma}

\begin{proof}
Every proxy remains in $S\setminus Q$ because $Q\subseteq S_0$.
The reference distances of good clients therefore sum to at most
$P_g+0.01P$.  For a covered bad cluster, its leaf gives the pointwise
bound
\begin{equation}
 d(p,H)\le d_p^*+r_i+R_i\qquad(p\in C_i).
 \label{sam:covered-client-bound}
\end{equation}
Use Equation~\eqref{sam:upward-radius}, $r_i\le1.2\bar d_i$, and
Lemma~\ref{sam:computable-radius} to obtain, with $P_A=\sum_{i\in A}P_i$,
\[
 \sum_{i\in A}\sum_{p\in C_i}d(p,H)
 \le(3.4+1.2\kappa)P_A+5.2\kappa P.
\]
For an uncovered nonempty cluster, retain its near-client cost $T_i$
and sum Equation~\eqref{sam:far-client-bound} over the others.  Their
number is at most $n_i$, so this cluster contributes at most
$1.85P_i+61T_i$.  Empty clusters contribute zero.  With
$P_B=\sum_{i\in B}P_i$ this yields
\begin{align*}
 W(H)
 &\le P_g+0.01P+(3.4+1.2\kappa)P_A+5.2\kappa P
                         +1.85P_B+61T_{\rm near}(H)\\
 &\le(3.41+6.4\kappa)P+61T_{\rm near}(H)
 <3.42P+61T_{\rm near}(H).
\end{align*}
Here $P_g+P_A+P_B=P$.

If $T=T_{\rm near}(H)>\kappa P$, the ball event has probability
\[
 \frac{T}{2W(H)}\ge\frac{T}{2(3.42P+61T)}\ge\rho_\kappa.
\]
If instead $T\le\kappa P$ and $I(H)>\kappa P$, then
$W(H)\le(3.42+61\kappa)P$.  The removal event has probability
\[
 \frac{I(H)}{2W(H)}>\frac{\kappa}{2(3.42+61\kappa)}=\rho_\kappa.
\]
In either nonterminal case $W(H)>0$, so
the client distribution is defined.
\end{proof}

\begin{theorem}[A successful sampling prefix]\label{sam:successful-prefix}
Assume the equal-size comparison setup and
Equation~\eqref{sam:proxy-hypothesis}, with $1\le m\le M_0$, and use
$\kappa=10^{-6}$.  Given $S$ and the correct guesses $P,m$, one run of the
actual experiment has, with probability at least $\exp(-C_\kappa m)$ for
an absolute constant $C_\kappa<\infty$, a recorded prefix with the
following properties:
\begin{enumerate}[label=\textup{(\roman*)}]
\item $Q\subseteq S_0$, and the balls correspond to distinct
      $i\in A\subseteq\mathcal B$, one per such cluster, with
      $w_i\in N_i$ and
      \[
       r_i=d(w_i,i)=d(w_i,O)\le R_i
       \le(1+\kappa)r_i+\kappa r^0(w_i).
      \]
      Consequently $i\in\mathcal F_i$ and
      $\sum_{i\in A}n_i r^0(w_i)\le5.2P$.
\item For $H$ in Equation~\eqref{sam:reference}, $B=\mathcal B\setminus A$,
      $R^*=S_0\setminus Q$, and the nearest assignment sets $J_s$,
      \[
       \sum_{i\in B}\sum_{p\in N_i}d(p,H)\le\kappa P,
       \qquad
       \sum_{s\in R^*}\sum_{p\in J_s}d(p,H)\le\kappa P.
      \]
\item The reference retains every good proxy.  The covered-client bound
      in Equation~\eqref{sam:covered-client-bound} holds, and for each
      uncovered nonempty cluster the far-client bound in
      Equation~\eqref{sam:far-client-bound} holds with
      $T_i=\sum_{p\in N_i}d(p,H)$.
\end{enumerate}
The run has at most $2m+1$ recorded prefixes and takes polynomial time
for fixed $\kappa$, with rational numbers of polynomial bit length.
Neither its input nor its stopping decisions require $O$, its partition,
the proxies, or $S_0$.
\end{theorem}

\begin{proof}
Define the guide as follows.  At a correct state satisfying
Equation~\eqref{sam:stopping-rule}, stop.  Otherwise, if the first
inequality fails, condition the actual client-and-mode draw on ball mode
and a client in $\bigcup_{i\in B}N_i$.  Such a client has a unique
assigned comparison cluster $i\in B$; add it to $A$ with the prescribed
correct radius and its leaf.  If only the second inequality fails,
condition on removal mode and a client assigned to $R^*$, and delete
that original center.  Lemma~\ref{sam:sampling-mass} supplies the same
fixed positive lower bound $\rho_\kappa$ on both conditioning events.
Every guided ball addition covers a new nonempty bad cluster, and every
guided removal deletes a new member of $S_0$.  Thus all its states remain
correct and all good proxies persist.

There are at most $m$ possible removals and at most $m$ distinct clusters
that can receive balls.  If all nonempty bad clusters have received balls
and all of $S_0$ has been removed, both stopping quantities are zero.
More generally, a nonterminal state always has one of the progress
operations just described.  Each such operation increases $|A|+|Q|$ by
one, a quantity at most $2m$.  The guide therefore stops within $2m$
steps.  Empty bad clusters can remain in $B$ throughout; they never
contribute to the stopping quantities or to any division by $n_i$.

Apply Lemma~\ref{sam:probability} with $\rho=\rho_\kappa$.  Its
exponential probability bound gives a recorded terminal prefix, and
the guide's invariants give all the stated properties, including
Lemma~\ref{sam:computable-radius} for the radius sum.  The arithmetic and
running-time assertions follow from the implementation below.
\end{proof}

\subsection{Enumeration, precision, and running time}

Here are full implementation details for the finite experiments above.
Let the largest original distance be the positive integer $\Delta$.
For every comparison solution $P$ is an integer satisfying
$1\le P\le n\Delta$.  Enumerating this whole range and
$m\in\{1,\ldots,\min\{\lfloor M_0\rfloor,h\}\}$ includes the
correct pair.  The $m=0$ candidate is simply $S$.  These ranges have
polynomial size on the integral instance, and the experiment for any
guessed pair has the same polynomial running-time bound, whether or not
that guess is correct.

Equation~\eqref{sam:r0-formula} computes $r^0(w)$ by sorting $n$ original
integer distances and taking the minimum of $n$ rational numbers.  Thus
it has polynomial bit length.  The first and third menus multiply it
by a rational $j\kappa$, with $j=O_\kappa(m+1)$.  The relative menu
uses a distance $d(w,w')+R'$ to an existing leaf and divides it by
$(1+\kappa)^e$, where $e\le E_\kappa$ is fixed independently of the
input size.  In a sequence of at most $2m$ steps, each such dependency
therefore increases numerator and denominator bit lengths by at most
the bit length of one original distance plus a constant depending on
$\kappa$.  All radii have polynomial bit length, including on incorrect
histories.  Distances to leaves, reference minima, client assignments,
the total weight, and facility-ball membership are consequently computable
in polynomial time with rational arithmetic.  Taking a common denominator
for a polynomial number of polynomial-bit weights still uses polynomially
many bits: the product of their denominators is already such a common
denominator, and its bit length is the sum of theirs.

For completeness, exact rational distributions need not be implemented by
an unbounded rejection loop.  At a state, enumerate the full finite
one-step outcomes: a client and removal mode, or a client, ball mode,
menu, and entry.  Include undefined choices as outcomes that end the run.
This is a polynomial number of outcomes for fixed $\kappa$, and their
exact rational probabilities $p_1,\ldots,p_t$ have polynomial bit length.
Choose an integer $b$ such that $2^{-b}\le p_j/2$ for every positive
$p_j$; polynomial bit length makes a polynomial value of $b$ sufficient
and computable.  Assign to outcome $j$ initially
$\lfloor2^bp_j\rfloor$ of the $2^b$ equiprobable binary strings, and
assign the leftover strings to any one outcome.  The resulting dyadic
probability is at least $p_j/2$ for every positive $p_j$.  Cumulative
integer counts implement the draw using exactly $b$ random bits and
polynomial arithmetic.  All actual transitions retain their specified
form, and Lemma~\ref{sam:probability} already accounts for the factor
$1/2$ lost at each transition.  This proves a polynomial worst-case
running-time bound with a bounded number of random bits.

Finally, the factor $\exp(O_\kappa(m))$ is polynomial when
$m\le M_0=O(\log N)$ and $\kappa$ is fixed.  Repeating the experiment
$\exp(O_\kappa(M_0))\log(1/\zeta)$ times makes the probability of
missing all successful prefixes at most $\zeta$.  Each run saves only
$2m+1$ states.  The completion procedure can therefore examine every
saved state, retain its cheapest feasible output together with the
fallback $S$, and amplify its own additional success probability without
ever recognizing the analytical stopping conditions.
